\documentclass[11pt]{article}
\usepackage[a4paper,margin=1in]{geometry}
\usepackage{amsmath,amssymb,amsthm}
\usepackage[T1]{fontenc}
\usepackage{lmodern}
\usepackage[hidelinks]{hyperref}
\setlength{\emergencystretch}{3em}

\theoremstyle{plain}
\newtheorem{theorem}{Theorem}
\newtheorem{lemma}[theorem]{Lemma}
\newtheorem{proposition}[theorem]{Proposition}
\newtheorem{corollary}[theorem]{Corollary}
\theoremstyle{remark}
\newtheorem{remark}[theorem]{Remark}
\newtheorem{example}[theorem]{Example}

\title{A Proof of Conjecture 00000000047:\\
Every $[n]$ Has a Prime Permutation}
\author{%
  Feiyu\thanks{Independent researcher.}
  \and
  Claude (Opus 5.5)\thanks{Anthropic.}%
}
\date{2026-10-02}

\begin{document}
\maketitle

\begin{abstract}
Conjecture 00000000047 asserts that for every sufficiently large even $n$
there is a permutation $\pi$ of $[n]=\{1,\ldots,n\}$ such that $\pi(i)+i$ is
prime for all $i$. We prove that such a permutation exists for every $n$, even
or odd, and that it can be chosen to be an involution. The proof is a short
induction on Bertrand's postulate, and it is formalized in Lean~4 against
Mathlib.
\end{abstract}

\section{The conjecture as stated}

The original text of conjecture 00000000047 reads:

\begin{quote}
Conjecture: For every sufficiently large even $n$ there exists a permutation
$\pi$ of $[n]$ such that $\pi(i)+i$ is prime for all $i$ (a ``prime
permutation'').
\end{quote}

\section{Reading of the statement}

As usual, $[n]=\{1,2,\ldots,n\}$. A \emph{prime permutation} of $[n]$ is a
bijection $\pi\colon[n]\to[n]$ such that $\pi(i)+i$ is a prime number for every
$i\in[n]$. The conjecture asserts that there is an $N$ such that every even
$n\ge N$ admits a prime permutation of $[n]$. There is no ambiguity to resolve.

\section{Result}

\begin{theorem}\label{thm:main}
For every integer $n\ge0$ there is an involution $\pi$ of $[n]$ such that
$\pi(i)+i$ is prime for all $i\in[n]$. In particular Conjecture 00000000047
holds, with any threshold $N$, and the evenness hypothesis is not needed.
\end{theorem}

\section{Proof}

We use Bertrand's postulate in the form: for every $n\ge1$ there is a prime $p$
with $n<p\le 2n$.

\begin{proof}[Proof of Theorem~\ref{thm:main}]
By strong induction on $n$. For $n=0$ the empty permutation works.

Let $n\ge1$, and assume the statement for all smaller values. Choose a prime
$p$ with $n<p\le2n$ and put $m=p-n$, so that $1\le m\le n$. Split
\[
  [n]=\{1,\ldots,m-1\}\ \cup\ \{m,\ldots,n\}.
\]
On the second block define $\pi(i)=p-i$. If $m\le i\le n$ then
$m=p-n\le p-i\le p-m=n$, so $\pi$ maps the block $\{m,\ldots,n\}$ into
itself; clearly $\pi(\pi(i))=i$, and $\pi(i)+i=p$ is prime. The first block is
$[m-1]$, and $m-1=p-n-1<n$ because $p\le 2n$. By the induction hypothesis there
is an involution of $[m-1]$ with the required property, and we let $\pi$ act on
the first block by it. The two pieces together form an involution of $[n]$, and
$\pi(i)+i$ is prime for every $i\in[n]$.
\end{proof}

\begin{example}
For $n=10$ take $p=11$, so $m=1$ and $\pi(i)=11-i$ on all of $[10]$: the pairs
$\{1,10\},\{2,9\},\{3,8\},\{4,7\},\{5,6\}$ all sum to $11$. For $n=8$ take
$p=11$, so $m=3$: the block $\{3,\ldots,8\}$ is paired as
$\{3,8\},\{4,7\},\{5,6\}$, and the remaining $[2]$ is handled with $p=3$,
giving the pair $\{1,2\}$. So $\pi=(1\,2)(3\,8)(4\,7)(5\,6)$ works for $n=8$.
\end{example}

\section{Formalization}

The accompanying Lean~4 project (\texttt{lean/Submission/Basic.lean}) states
the conjecture and proves it against Mathlib.

\begin{itemize}
  \item \texttt{HasPrimePermutation n} --- there is
    \texttt{$\pi$ : Equiv.Perm (Finset.Icc 1 n)} with
    \texttt{Nat.Prime ($\pi$ i + i)} for every \texttt{i}. A permutation of the
    finite set $\{1,\ldots,n\}$, as a type, is exactly a permutation of $[n]$.
  \item \texttt{ConjectureHolds} --- for some $N$, every even $n\ge N$
    satisfies \texttt{HasPrimePermutation n}.
  \item \texttt{exists\_prime\_involution} --- the induction above, producing
    a function $f\colon\mathbb N\to\mathbb N$ that maps $[n]$ into itself,
    satisfies $f(f(i))=i$ and makes $f(i)+i$ prime for $i\in[n]$. It uses
    Mathlib's Bertrand postulate
    \texttt{Nat.exists\_prime\_lt\_and\_le\_two\_mul}.
  \item \texttt{hasPrimePermutation} --- every $n$ admits a prime permutation,
    obtained from $f$ via \texttt{Function.Involutive.toPerm}.
  \item \texttt{conjecture\_00000000047} --- \texttt{ConjectureHolds}, with
    threshold $N=0$.
\end{itemize}

Primality is Mathlib's \texttt{Nat.Prime}. The toolchain is
\texttt{leanprover/lean4:v4.33.1}, Mathlib is pinned to \texttt{v4.33.1}, and
\texttt{lake build} succeeds. The project contains no \texttt{sorry}, no \texttt{admit} and no
\texttt{native\_decide}, and introduces no axioms:
\begin{center}
\texttt{\#print axioms conjecture\_00000000047}
\end{center}
reports exactly \texttt{propext}, \texttt{Classical.choice} and
\texttt{Quot.sound}.

\section{Remarks}

\begin{remark}[fixed points]
For even $n$ the involution constructed above has no fixed points, so it
splits $[n]$ into $n/2$ pairs with prime sums. Indeed a fixed point needs
$p=2i$, so $p=2$, which happens only in the case $n=1$. When $n$ is even, the
paired block $\{m,\ldots,n\}$ has $2n-p+1$ elements, an even number because $p$
is odd. So the leftover block $[m-1]$ again has an even number of elements, and
the case $n=1$ never arises. This pairing statement for even $n$ is a classical
consequence of Bertrand's postulate; see Greenfield and Greenfield
\cite{greenfield}, who also discuss harder variants. For odd $n$ a fixed point
is unavoidable for an involution, and our construction places it at $i=1$,
where $1+1=2$ is prime. This fixed-point refinement is not part of the Lean
development.
\end{remark}

\begin{remark}
The conjecture asks only for large even $n$, which suggests a probabilistic
argument was intended. Nothing of the kind is needed: the deterministic
construction covers all $n$. A related but genuinely different question is
whether the prime-sum graph on $[n]$ has a \emph{Hamilton cycle}
(Conjecture 00000000045). There the vertices must form a single cycle, not a
union of transpositions, and parity already rules it out for odd $n$.
\end{remark}

\begin{thebibliography}{9}
\bibitem{greenfield} L.~E.~Greenfield and S.~J.~Greenfield, \emph{Some
  problems of combinatorial number theory related to Bertrand's postulate},
  J. Integer Sequences 1 (1998).
\end{thebibliography}

\end{document}
